%!TEX root = ../template.tex
%% ====================================================================
%% Capitulo 8 -- Conclusoes e trabalho futuro
%% Esqueleto gerado 2026-08-13 a partir do indice v1
%% (Mind Palace: salas/tese/2026-08-13-indice-dissertacao-v1.md)
%% Cada '>>' descreve o que a seccao deve conter. Apagar ao escrever.
%% ====================================================================

\chapter{Conclusões e trabalho futuro}
\label{cha:conclusoes}

\section{Conclusões}
\label{sec:conclusoes}
% [REV-2026-09-13][C8-01][P2][CONTEUDO]
\revisaotese{C8-01}{P2}{CONTEUDO}
{Fechar perguntas com resultados e limites, sem novo manifesto}
{Escrever uma resposta curta por pergunta/objetivo alinhado no Cap. 1:
plataforma e validação efetivamente verificadas; diferenças entre
vetores e sensibilidade; decisão dos candidatos sob o critério adotado;
contributo/estado da revisão; deteção real não estimada.
Usar 3--5 resultados com denominadores já auditados. Não repetir que
as normas só exigem a primeira camada, que zero candidatos prova
impossibilidade da classe ou que foram demonstradas poupanças.}%
% [/REV-2026-09-13]

%% >> Fiabilidade exige dados reconciliados E modelos validados; as normas so obrigam a primeira metade.


\section{Roteiro de modelos sucessores}
\label{sec:roteiro-modelos}
% [REV-2026-09-13][C8-02][P2][ESTRUTURA]
\revisaotese{C8-02}{P2}{ESTRUTURA}
{Trocar escada de complexidade por decisões guiadas pelo diagnóstico}
{ARMA, estado, GAM, pooling e grey-box não formam uma sequência em que
mais complexo é melhor. Recomenda-se fundir esta secção com Trabalho
futuro numa tabela: falha observada -\textgreater{} opção -\textgreater{} dados adicionais -\textgreater{}
comparador -\textgreater{} critério de sucesso. Primeiro calendário/eventos,
fronteiras e referência verificáveis; depois piloto prospetivo.
Uma referência constante bem avaliada deve continuar como comparador
dos vapores, e o fuel gás como candidato a piloto, sem prometer ganho.}%
% [/REV-2026-09-13]

%% >> Escada fundamentada pelo diagnostico: erros ARMA -> espaco de estados/Kalman (deriva como estado estimado) -> GAM -> hierarquico multi-unidade -> grey-box fisico.


\section{Trabalho futuro}
\label{sec:trabalho-futuro}

%% >> Piloto de modelo dinamico no fuel gas, expansao com pooling parcial, analise custo/beneficio.


